\section{Forschungsfragen}
Ausgehend von dem aktuellen Forschungsstand zeigt sich, dass diese zwar umfangreiche Erkenntnisse über die Nutzung und Entwicklung des Online-Shoppings während der Corona-Pandemie liefern, diese aber nur begrenzt einen Einblick in die subjektiven Bedeutungszuschreibungen sowie Erfahrungen von Konsumenten geben. Vor diesem Hintergund wurde die folgende Forschungsfrage formuliert:

\enquote{Wie hat die Corona-Pandemie das individuelle Kaufverhalten zwischen Online-Shopping und stationärem Einzelhandel verändert?}

Neben der Hauptforschungsfrage wurden auch noch weitere Unterfragen formuliert:

\begin{itemize}
    \item Welche Motive führten während der Corona-Pandemie zu einer verstärkten Nutzung des Online-Shoppings?
    \item Wie bewerten Konsumenten das Zusammenspiel von Online-Shoppinh und dem stationären Handel?
    \item Welche Faktoren haben Einfluss auf die Entscheidung der Konsumenten, zwischen Online-Shopping und dem stationärem Einzelhandel zu wechseln?
\end{itemize}

Diese Fragen wurden offen formuliert, um eine explorative Unterscuhung mittels qualitativer Forschungsmethoden zu ermöglichen.